\begin{agradecimentos}
	Primeiramente, a Deus, pela força, resiliência e clareza mental concedidas ao longo de toda esta jornada, permitindo-me superar os momentos de cansaço e os desafios inerentes à pesquisa e ao desenvolvimento de software.
	
	Ao Instituto Federal de Educação, Ciência e Tecnologia do Tocantins (IFTO) - Campus Colinas do Tocantins, por ser o alicerce da minha formação profissional. Aos excelentes professores do curso de Tecnologia em Análise e Desenvolvimento de Sistemas (TADS), que não apenas compartilharam seus conhecimentos teóricos e práticos, mas também me inspiraram a buscar sempre a excelência e a inovação tecnológica.
	
	Um agradecimento muito especial ao meu orientador, Prof. Me. Fernando Turibio de Moura. Sua paciência, dedicação e vasto conhecimento foram fundamentais para transformar uma ideia ambiciosa em um projeto real. Obrigado pela confiança no meu potencial, pelas revisões minuciosas e por me guiar com maestria pelas melhores práticas de engenharia de software, especialmente nos momentos mais complexos de integração deste simulador.
	
	À minha esposa e à minha família, pela compreensão inesgotável nas inúmeras noites, madrugadas e finais de semana dedicados à programação, à correção de bugs e à escrita deste documento. Sem o suporte emocional, o incentivo constante e a paciência de vocês com as minhas ausências, a conclusão desta etapa não seria possível. Vocês são a minha maior motivação.
	
	À gigantesca comunidade de desenvolvedores de software livre, em especial aos criadores e mantenedores da Godot Engine, da linguagem Python e do motor de xadrez Stockfish. A filosofia de código aberto, a documentação acessível e o compartilhamento global de conhecimento foram a base técnica que viabilizou a materialização e o sucesso deste sistema.
	
	Aos meus colegas de turma e amigos de curso, por todas as linhas de código revisadas em conjunto, pelos testes incansáveis do jogo em suas diversas versões, pelas partidas jogadas e pelas horas de estudo compartilhadas. Vocês tornaram os desafios desta graduação muito mais leves e a jornada muito mais inesquecível.
	
	Por fim, a todos que, de forma direta ou indireta, contribuíram com palavras de incentivo, críticas construtivas e apoio ao longo destes anos de graduação, meu mais sincero muito obrigado.
\end{agradecimentos}